% !TEX root = ../main.tex

\chapter{Experimental Setup}\label{chapter:experiments}

\section{Dataset and Splits}\label{sec:dataset}

The dataset is a MATSim simulation of the Paris metropolitan area provided by Natterer~et~al.~\cite{natterer2025ml}.\ Each scenario corresponds to a distinct policy intervention: a random subset of road segments is selected and assigned a capacity reduction in $\{10\%, 20\%, \dots, 100\%\}$.\ MATSim is then run from the same baseline plans, and the per-segment difference in hourly volume is recorded as the supervised target $\Delta v$ (Table~\ref{tab:dataset}).

\begin{table}[htbp]
 \centering
 \caption[Paris MATSim dataset statistics]{Paris MATSim dataset.\ All experiments use the 1{,}000-scenario subset.}
 \label{tab:dataset}
 \begin{tabular}{lr}
 \toprule
 Property & Value \\
 \midrule
 Road segments per graph & 31{,}635 \\
 Available scenarios (full) & 10{,}000 \\
 Scenarios used (10\% subset) & 1{,}000 \\
 Test scenarios (T1--T6 split 80/15/5) & 50 \\
 Test scenarios (T7--T11, DE; split 80/10/10) & 100 \\
 Test predictions, T7--T11, DE & 3{,}163{,}500 \\
 \bottomrule
 \end{tabular}
\end{table}

Trials~T1--T6 use an 80/15/5 split (800/150/50 scenarios); Trials~T7--T11 and the Deep Ensemble use 80/10/10 (800/100/100), with the larger 100-graph test set enabling the high-resolution UQ analyses in Chapter~\ref{chapter:results}.\ The split is at the scenario level (each scenario appears in exactly one split) with split seed $42$, ensuring identical scenario assignment across all trials that share a split configuration.

\paragraph{A note on the 10\% subset.}\ The full Natterer corpus is around 200 batch files of 50 graphs each ($\approx 10{,}000$ scenarios); we load the first 20 batches, which gives the $1{,}000$-scenario subset reused across all eleven trials so that method-to-method differences are evaluated on identical data.\ The 100-graph evaluation set ($3{,}163{,}500$ node-level predictions per trial) is large enough for stable per-graph statistics: per-graph Spearman~$\rho$ has std $0.023$, with a bootstrap 95\% CI of $\pm 0.005$ on the mean.\ Possible subset-selection bias from the contiguous batch ordering, and the question of whether the findings survive at the full $10{,}000$-scenario scale, are discussed under threats to validity in Section~\ref{sec:threats_to_validity}.

\section{Training Protocol}\label{sec:training_protocol}

Eight base trials of the PointNetTransfGAT architecture are evaluated.\ Trial~1 uses a \texttt{Linear} final layer; T2--T8 use \texttt{GATConv}.\ Hyperparameters are listed in Table~\ref{tab:trials}.

\begin{table}[htbp]
 \centering
 \caption[Training trials and hyperparameters]{Training trials.\ T1 uses a Linear final layer; T2--T8 use GATConv.}
 \label{tab:trials}
 \begin{tabular}{lcccccc}
 \toprule
 Trial & Batch & Split & Effective dropout & LR & Weighted MSE & Final layer \\
 \midrule
 T1 & 32 & 80/15/5 & 0.0$^\dagger$ & $10^{-3}$ & No & Linear \\
 T2 & 16 & 80/15/5 & 0.3 & $5 \times 10^{-4}$ & No & GATConv \\
 T3 & 16 & 80/15/5 & 0.0 & $5 \times 10^{-4}$ & Yes & GATConv \\
 T4 & 16 & 80/15/5 & 0.0 & $5 \times 10^{-4}$ & Yes & GATConv \\
 T5 & 8 & 80/15/5 & 0.3 & $5 \times 10^{-4}$ & No & GATConv \\
 T6 & 8 & 80/15/5 & 0.3 & $3 \times 10^{-4}$ & No & GATConv \\
 T7 & 8 & 80/10/10 & 0.3 & $6 \times 10^{-4}$ & No & GATConv \\
 \textbf{T8} & \textbf{8} & \textbf{80/10/10} & \textbf{0.2} & $\mathbf{5 \times 10^{-4}}$ & \textbf{No} & \textbf{GATConv} \\
 \bottomrule
 \end{tabular}
 \\[2pt]
 \footnotesize $^\dagger$ T1 was instantiated with \texttt{use\_dropout=False}, so the effective dropout rate is $0$; MC Dropout is therefore undefined on T1.\ T3 and T4 likewise have effective dropout zero and are not MC-Dropout-compatible.
\end{table}

All trials use the AdamW optimiser~\cite{loshchilov2019decoupled} with weight decay $10^{-4}$ and MSE loss.\ T3 and T4 use a weighted variant $\mathcal{L}_\text{weighted} = \tfrac{1}{n}\sum_i w_i(y_i - \hat y_i)^2$, with weights $w_i$ proportional to $|y_i|$ so that non-zero-$\Delta v$ examples carry more influence.\ T3 and T4 share identical hyperparameters; T4 is a replication run of T3 with a different random seed, included to check that the weighted-MSE failure mode is systematic rather than a one-off.\ The implementation uses PyTorch~\cite{paszke2019pytorch} and PyTorch Geometric~\cite{fey2019fast}; hardware is described in Section~\ref{sec:resources}.\ The early-stopping criterion is validation $R^2$ with a patience of 25 epochs and a minimum improvement of $10^{-3}$.\ Among Trials~2--8, T8 attains the highest test $R^2 = 0.5957$ on the 100-graph evaluation set; together with active dropout ($p = 0.2$), this makes it the strongest UQ-compatible base in this exploratory comparison, and it is used as the primary model for the UQ analyses that follow.

\section{Post-Hoc UQ and Ensemble Experimental Design}\label{sec:uq_design}

Post-hoc analyses on Trial~8 and the ensemble experiments are summarised in Table~\ref{tab:exp_design}.\ All operate on the same 100-graph T8 test set ($3{,}163{,}500$ nodes) unless noted.

\begin{table}[htbp]
\centering
\caption[Post-hoc and ensemble experimental design]{Post-hoc UQ and ensemble experimental design.\ Scenario-level splits use seed $42$ throughout.}
\label{tab:exp_design}
\footnotesize
\setlength{\tabcolsep}{4pt}
\renewcommand{\arraystretch}{1.15}
\begin{tabular}{p{3.7cm}p{8.2cm}}
\toprule
\textbf{Experiment} & \textbf{Protocol} \\
\midrule
MC Dropout (primary)\label{sec:mc_design} & $S = 30$ stochastic forward passes; primary metric Spearman $\rho$ between $\sigma$ and $|y - \hat y|$. \\
$S$-convergence\label{sec:s_conv_design} & 10 test graphs, $S \in \{5, 10, \dots, 50\}$. \\
Bootstrap CI\label{sec:bootstrap_design} & Per-graph $\rho$ resampled $B = 10{,}000$ times for 95\% CI on the mean. \\
T7 replication\label{sec:t7_design} & Same MC Dropout / AUROC / conformal protocol on Trial~7. \\
Split conformal\label{sec:conformal_design} & Scenario-level 50/50 partition; nonconformity $|y - \hat y|$; coverage at $\{0.5, 0.7, 0.8, 0.9, 0.95\}$.\ Adaptive variant normalises by MC Dropout $\sigma$. \\
Calibration audit\label{sec:calibration_design} & Empirical $k_p$ at $p \in \{0.5, 0.7, 0.8, 0.9, 0.95\}$ on the full test set. \\
$\sigma$-scaling\label{sec:ts_design} & Fitted on a $30\%$ random node-level subset ($949{,}050$ nodes), evaluated on the remaining $70\%$; grid $T \in [0.5, 5.0]$, ECE-minimising. \\
Scoring + selective + AUROC\label{sec:scoring_design} & CRPS, PIT, Winkler on the full test set; selective prediction at $\tau \in \{0.10, 0.25, 0.50, 0.75, 1.00\}$; AUROC/AUPRC at top-$10\%$ / top-$20\%$. \\
\midrule
Experiment A (seed ensemble)\label{sec:expA_design} & Same trained T8, 5 inference seeds; compares $\sigma_\text{MC}$, $\sigma_\text{seed-ens}$, and the quadrature combination. \\
Experiment B (multi-model)\label{sec:expB_design} & $R^2$-weighted ensemble of T2/T5/T6/T7/T8, deterministic inference. \\
Trial D (Deep Ensemble)\label{sec:de_design} & 5 PointNetTransfGAT trained from scratch with seeds $\{42, 137, 256, 389, 512\}$, T8 hyperparameters, 80/10/10 split. \\
\bottomrule
\end{tabular}
\end{table}

\section{Uncertainty-Aware Training Extensions}\label{sec:uat_design}

Three trials evaluate whether uncertainty-aware training objectives can produce natively calibrated intervals while preserving T8's accuracy.\ All replace T8's final \texttt{GATConv(64$\to$1)} with a $134$-parameter \texttt{GATConv(64$\to$2)} head, differing in head interpretation, loss, and backbone trainability (Table~\ref{tab:uat_design}).\ Other hyperparameters (optimiser AdamW, learning rate $5\times10^{-4}$, batch 8, 80/10/10 split) inherit from T8.\ T9 freezes the backbone weights but keeps dropout active for MC Dropout epistemic estimation ($\sigma_\text{tot} = \sqrt{\sigma_\text{alea}^2 + \sigma_\text{epi}^2}$); T10/T11 use a single deterministic forward pass.

\begin{table}[htbp]
\centering
\caption[UAT design summary]{Uncertainty-aware training extensions.\ Head output channels are $(\hat\mu,\,\log\hat\sigma^2)$ for T9 and $(\hat q^\text{lo}, \hat q^\text{hi})$ for T10/T11 ($\tau = 0.05, 0.95$).}
\label{tab:uat_design}
\footnotesize
\setlength{\tabcolsep}{4pt}
\renewcommand{\arraystretch}{1.15}
\begin{tabular}{p{0.9cm}p{2.4cm}p{2.0cm}p{1.5cm}p{4.8cm}}
\toprule
Trial & Loss & Backbone & UQ output & Gates \\
\midrule
T9\label{sec:t9_design}  & Hetero NLL ($\lambda = 0.1$)~\cite{seitzer2022pitfalls} & Frozen & MC + $\hat\sigma_\text{alea}$ & $R^2 \geq 0.55$, $\mathrm{PICP}_{90} \geq 85\%$, $\mathrm{PICP}_{95} \geq 90\%$ \\
T10\label{sec:t10_design} & Pinball & Unfrozen & CQR + $\hat Q$ & $R^2 \geq 0.57$, $\mathrm{PICP}_{90} \geq 88\%$, $\mathrm{PICP}_{95} \geq 93\%$, $\hat Q > 0$, $W_{90} < W_{95}$ \\
T11\label{sec:t11_design} & Pinball & Frozen & CQR + $\hat Q$ & same as T10 \\
\bottomrule
\end{tabular}
\end{table}

The single design difference between T10 and T11 is the backbone-trainability flag; head, loss, calibration protocol and learning rate are identical, so the T10/T11 contrast isolates that flag.

\section{Computational Resources}\label{sec:resources}

All training and inference used a single NVIDIA T4 GPU (16~GB VRAM).\ Wall-clock runtimes from the saved \texttt{training\_log.json} files: the eight base trials totalled $\approx 13$~h; T8 MC Dropout inference ($S=30$, 100 graphs) $228$~min; Trials~9, 10 and 11 took $14.5$, $87$ and $39.7$~h respectively; the Deep Ensemble $\approx 5\times$ a single T8 run.\ Trial~10 is the longest because it fine-tunes the full $1{,}416{,}768$-parameter backbone end-to-end, while Trials~9 and 11 train only a $134$-parameter head with the backbone frozen.\ Non-GNN baselines (Section~\ref{sec:non_gnn_baselines}) were run separately on an A100 (80~GB) runtime.
